
\documentclass[12pt,a4paper]{report}

% Packages
\usepackage[a4paper, margin=1in]{geometry}
\usepackage{graphicx}
\usepackage{amsmath}
\usepackage{amssymb}
\usepackage{booktabs}
\usepackage{array}
\usepackage{tabularx}
\usepackage{float}
\usepackage{fancyhdr}
\usepackage[hidelinks]{hyperref}
\usepackage{xcolor}
\usepackage{setspace}
\usepackage{enumitem}
\usepackage{caption}

\captionsetup{
    labelfont=bf,
    justification=raggedright,
    singlelinecheck=false,
    format=plain
}
\usepackage{subcaption}
\usepackage{multirow}
\usepackage{colortbl}
\usepackage{longtable}
\setlength{\LTleft}{\fill}
\setlength{\LTright}{\fill}
\usepackage{titlesec}
\usepackage{titletoc}
\usepackage{lastpage}
\usepackage{makecell}
\usepackage{etoolbox}
\usepackage{needspace}

\numberwithin{table}{chapter}
\numberwithin{figure}{chapter}
% Table layout and spacing: consistent padding, row height, and longtable pre/post skips
\setlength{\tabcolsep}{6pt}
\renewcommand{\arraystretch}{1.12}
\setlength{\LTpre}{0pt}
\setlength{\LTpost}{6pt}
% Table rules (outline thickness) and small extra row height for clarity
\setlength{\arrayrulewidth}{0.5pt}
\setlength{\extrarowheight}{0.6pt}

% Prevent tables from overflowing past the page bottom:
% if fewer than 5 baseline-skips remain, break to the next page first.
\BeforeBeginEnvironment{table}{\needspace{5\baselineskip}}
\BeforeBeginEnvironment{longtable}{\needspace{5\baselineskip}}

\definecolor{osdagGreen}{HTML}{91B014}

\fancypagestyle{main}{
  \fancyhf{}
  \fancyhead[L]{PROJECTNEWB $|$ }
  \fancyhead[R]{2026-11-21 $|$ Rev 0}
  \fancyfoot[L]{Osdag $|$ FOSSEE $|$ Indian Institute of Technology Bombay}
  \fancyfoot[R]{Page \thepage\ of \pageref{LastPage}}
  \renewcommand{\headrule}{\color{osdagGreen}\hrule width\headwidth height 1pt \vspace{2pt}}
  \renewcommand{\footrule}{%
    \ifbool{hasSDonPage}{%
      \vspace{-20pt}%
      \hbox to \headwidth{\textcolor{black}{\footnotesize\textit{* Software default value}}\hfil}%
      \vspace{4pt}%
    }{%
      \vspace{-8pt}%
    }%
    \color{osdagGreen}\hrule width\headwidth height 1pt \vspace{6pt}%
  }
}
\fancypagestyle{plain}{
  \fancyhf{}
  \fancyhead[L]{PROJECTNEWB $|$ }
  \fancyhead[R]{2026-11-21 $|$ Rev 0}
  \fancyfoot[L]{Osdag $|$ FOSSEE $|$ Indian Institute of Technology Bombay}
  \fancyfoot[R]{Page \thepage\ of \pageref{LastPage}}
  \renewcommand{\headrule}{\color{osdagGreen}\hrule width\headwidth height 1pt \vspace{2pt}}
  \renewcommand{\footrule}{%
    \ifbool{hasSDonPage}{%
      \vspace{-20pt}%
      \hbox to \headwidth{\textcolor{black}{\footnotesize\textit{* Software default value}}\hfil}%
      \vspace{4pt}%
    }{%
      \vspace{-8pt}%
    }%
    \color{osdagGreen}\hrule width\headwidth height 1pt \vspace{6pt}%
  }
}
\fancypagestyle{firstpage}{
  \fancyhf{}
  \renewcommand{\headrulewidth}{0pt}
  \fancyfoot[L]{Osdag $|$ FOSSEE $|$ Indian Institute of Technology Bombay}
  \fancyfoot[R]{Page \thepage\ of \pageref{LastPage}}
  \renewcommand{\footrule}{\vspace{-8pt}\color{osdagGreen}\hrule width\headwidth height 1pt \vspace{6pt}}
}
\pagestyle{main}
\setstretch{1.15}

% Custom Commands
\newcommand{\placeholder}[1]{\textit{\textless #1\textgreater}}
\newcommand{\todo}[1]{\colorbox{yellow}{TODO: #1}}
\newcolumntype{L}[1]{>{\raggedright\arraybackslash}p{#1}}
\newcolumntype{C}[1]{>{\centering\arraybackslash}p{#1}}
\newcolumntype{R}[1]{>{\raggedleft\arraybackslash}p{#1}}

% Software-default asterisk
\newcommand{\sdstar}{\textsuperscript{*}}
\newbool{hasSDonPage}
\boolfalse{hasSDonPage}
\newcommand{\markSD}{\global\booltrue{hasSDonPage}}
\renewcommand{\sdstar}{\textsuperscript{*}\markSD{}}
\AddToHook{shipout/before}{\global\boolfalse{hasSDonPage}}

\title{\Large\textbf{OsdagBridge} \\ \normalsize Open Source Software for Steel Girder Bridge Design \\ \vspace{2cm} \large Design Report}
\author{}
\date{}

\begin{document}


\begin{titlepage}
\thispagestyle{firstpage}
\centering
\vspace*{1.5cm}
\noindent
\begin{minipage}[c]{0.6\textwidth}
\raggedright
\includegraphics[width=\linewidth,keepaspectratio]{assets/osdag_logo.png}
\end{minipage}%
\hfill
\begin{minipage}[c]{0.35\textwidth}
\raggedleft

\end{minipage}
\\[1cm]

{\Huge \textbf{OsdagBridge}}\\[0.3cm]
{\large Open Source Software for Steel Girder Bridge Design}\\[1.5cm]
{\Large Design Report}\\[1.5cm]
\begin{tabular}{|L{4cm}|L{10cm}|}
\hline
\textbf{Project Name} & PROJECTNEWB \\
\hline
\textbf{Project Location} & Ahmedabad \\
\hline
\textbf{Author / Designer} & QA \\
\hline
\textbf{Reviewer} &  \\
\hline
\textbf{Organization} & QA \\
\hline
\textbf{Client Name and Organization} & QA \\
\hline
\textbf{Job Number} &  \\
\hline
\textbf{Date} & 2026-11-21 \\
\hline
\textbf{Report Version} &  \\
\hline
\end{tabular}
\end{titlepage}


\newpage
{\centering\Large\bfseries Executive Summary\par}
\addcontentsline{toc}{chapter}{Executive Summary}
\vspace{0.8em}

This section provides a concise summary of the bridge design, key inputs, governing loads, and final design outcomes.

\section*{Project Overview}
\addcontentsline{toc}{section}{Project Overview}
\label{sec:project-overview}


\begin{tabular}{|L{5.5cm}|L{8.5cm}|}
\hline
\textbf{Bridge Type} & Highway Bridge \\
\hline
\textbf{Design Standard} & IRC 5, IRC 6, IRC 22, IRC 24, IS 800 \\
\hline
\textbf{Span} & 30.0 m \\
\hline
\textbf{Carriageway Width} & 10.5 m \\
\hline
\textbf{No. of Girders} & 5 \\
\hline
\textbf{Girder Spacing} & 2.28 \\
\hline
\textbf{Deck Thickness} & 250.0 \\
\hline
\textbf{Overall Design Status} & Fail (Fatigue Normal Stress, Int.Stiff:\allowbreak{} Leg ≤ H\_\allowbreak{}limit, Fatigue Shear Stress) \\
\hline
\textbf{Governing Check} & Fatigue Normal Stress \\
\hline
\textbf{Overall Utilization Ratio (max)} & 1.87 (Girder) \\
\hline
\end{tabular}


\begin{figure}[H]
\centering
\fbox{\parbox{0.97\textwidth}{
\textit{[ PLACEHOLDER: Figure 1 -- Overall Bridge Plan ]}
}}
\caption*{Figure 1 -- Overall Bridge Plan}
\end{figure}

\newpage

\begin{figure}[H]
\centering
\fbox{\parbox{0.97\textwidth}{
\textit{[ PLACEHOLDER: Figure 2 -- Typical Cross-Section (with girder, deck, barriers, footpath) ]}
}}
\caption*{Figure 2 -- Typical Cross-Section (with girder, deck, barriers, footpath)}
\end{figure}

\begin{figure}[H]
\centering
\fbox{\parbox{0.97\textwidth}{
\textit{[ PLACEHOLDER: Figure 3 -- 3D View of Bridge Superstructure ]}
}}
\caption*{Figure 3 -- 3D View of Bridge Superstructure}
\end{figure}

\noindent\textbf{Table 1 -- Final Bridge Geometry (after optimization)}

\vspace{0.4em}
\noindent
\begin{tabular}{|C{2.8cm}|C{2.4cm}|C{2.4cm}|C{2.4cm}|C{2.4cm}|C{2.4cm}|}
\hline
  \textbf{} &
  \textbf{G1} &
  \textbf{G2} &
  \textbf{G3} &
  \textbf{G4} &
  \textbf{G5} \\
\hline
Member ID & G1M1 & G2M1 & G3M1 & G4M1 & G5M1 \\
\hline
Section Designation & 1670 x 510 x 22 x 510 x 22 & 1670 x 510 x 22 x 510 x 22 & 1670 x 510 x 22 x 510 x 22 & 1670 x 510 x 22 x 510 x 22 & 1670 x 510 x 22 x 510 x 22 \\
\hline
Governing Check & Fatigue Normal Stress & Fatigue Normal Stress & Fatigue Normal Stress & Fatigue Normal Stress & Fatigue Normal Stress \\
\hline
Utilization Ratio & 1.87 (Girder) & 1.87 (Girder) & 1.87 (Girder) & 1.87 (Girder) & 1.87 (Girder) \\
\hline
\end{tabular}

\vspace{0.4em}
\noindent\textit{Note: Utilization ratio (UR) = demand / capacity. A value $< 1.0$ indicates a passing check.}

\vspace{1em}

\section*{Key Design Outcomes Summary}
\addcontentsline{toc}{section}{Key Design Outcomes Summary}
\label{sec:key-outcomes}

\noindent Girder design pass \\
Cross bracing design pass \\
End Diaphragm design pass \\
Deck design pass

\section*{Design Assumptions and Limitations}
\addcontentsline{toc}{section}{Design Assumptions and Limitations}
\label{sec:assumptions}

\begin{itemize}
\item Additional inputs not provided by the user were assumed by software per IRC/IS code defaults or practical consideration.
\item Grillage analysis was performed using OSPGrillage assuming simply supported I-girders.
\item Substructure and foundation design are not included in this report.
\item Splice connections and bearings are not designed in this version.
\end{itemize}

% Restore numbered chapter format
\titleformat{\chapter}[block]{\normalfont\Large\bfseries\centering}{\thechapter}{1em}{}
\titlespacing*{\chapter}{0pt}{-30pt}{10pt}


\chapter{Project Information}

This section records all project metadata as entered by the designer.

\section{Project and Design Team Details}
\label{sec:project-details}

\begin{tabular}{|L{5.5cm}|L{8.5cm}|}
\hline
\textbf{Project Name} & PROJECTNEWB \\
\hline
\textbf{Project Location} & Ahmedabad \\
\hline
\textbf{Designer} & QA \\
\hline
\textbf{Reviewer} &  \\
\hline
\textbf{Organization} & QA \\
\hline
\textbf{Client} & QA \\
\hline
\textbf{Software Version} & OsdagBridge \\
\hline
\end{tabular}


\section{Applicable Codes and Standards}
\label{sec:codes}

\begin{itemize}
\item Indian Roads Congress (IRC) 5: General Features of Design
\item Indian Roads Congress (IRC) 6: Loads and Load Combinations
\item Indian Roads Congress (IRC) 22: Composite Construction (Limit State Design)
\item Indian Roads Congress (IRC) 24: Steel Road Bridges (Limit State Method)
\item Indian Roads Congress (IRC) 112: Concrete Road Bridges (deck design)
\item Indian Roads Congress Special Publication (IRC SP) 114: Seismic Design of Road Bridges
\item Indian Standard (IS) 800: General Construction in Steel
\item Indian Standard (IS) 2062: Hot Rolled Structural Steel Specification
\item Indian Standard (IS) 6006: Steel Bearings
\end{itemize}


\chapter{Input Parameters}

\setlength{\abovecaptionskip}{2pt}
\setlength{\belowcaptionskip}{2pt}

This section documents all inputs provided to OsdagBridge. User-provided inputs are clearly distinguished from software-assumed defaults. Where the user did not supply a value, the software has applied the IRC/IS code default or an empirical guideline; these are annotated with an asterisk (\sdstar{}).

\section{Basic Inputs (User-Defined)}
\label{sec:basic-inputs}

\noindent\textit{Note: These inputs are mandatory and were provided by the user.}

\begin{table}[H]
\caption{\textbf{Project Location}}
\label{subsec:project-location}
\begin{tabular}{|L{5.5cm}|L{8.5cm}|}
\hline
\textbf{parameter} & \textbf{value} \\
\hline
\textnormal{Project Location} & Ahmedabad \\
\hline
\textnormal{Latitude / Longitude} & ,  \\
\hline
\textnormal{Seismic Zone (IRC 6)} &  \\
\hline
\textnormal{Basic Wind Speed (IRC 6)} &  \\
\hline
\textnormal{Shade Temp. Max / Min (IRC 6)} &  °C /  °C \\
\hline
\end{tabular}
\vspace{0.4cm}
\end{table}

\begin{table}[H]
\caption{\textbf{Bridge Geometry}}
\label{subsec:bridge-geometry}
\begin{tabular}{|L{5.5cm}|L{8.5cm}|}
\hline
\textbf{parameter} & \textbf{value} \\
\hline
\textnormal{Type of Structure} & Highway Bridge \\
\hline
\textnormal{Span (m)} & 30.0 m \\
\hline
\textnormal{Carriageway Width (m)} & 10.5 m \\
\hline
\textnormal{Include Median} & No \\
\hline
\textnormal{Footpath} & None \\
\hline
\textnormal{Skew Angle (degrees)} & 0.0° (IRC 24 Cl. 504.8 limit: $\pm$15°) \\
\hline
\end{tabular}
\end{table}
\vspace{0.4cm}

\begin{table}[H]
\caption{\textbf{Material Selection}}
\label{subsec:material}
\begin{tabular}{|L{5.5cm}|L{8.5cm}|}
\hline
\textbf{parameter} & \textbf{value} \\
\hline
\textnormal{Girder Steel Grade (IS 2062)} & E 350A \\
\hline
\textnormal{Cross Bracing Steel Grade} & E 350A \\
\hline
\textnormal{End Diaphragm Steel Grade} & E 350A \\
\hline
\textnormal{Concrete Deck Grade (IRC 22)} & M40 \\
\hline
\end{tabular}
\end{table}
\vspace{0.4cm}

\newpage
\section{Additional Inputs}
\label{sec:additional-inputs}

Where the user has modified additional inputs, those values are reported here. Where no modification was made, the software default is shown.

\vspace{0.8cm}

\begin{longtable}{|L{5.5cm}|p{10.0cm}|}
\caption{\textbf{Typical Section Details}}
\hline
\textbf{parameter} & \textbf{value} \\
\hline
\textnormal{Overall Bridge Width (m)} & 11.399999999999999 \\[6pt]
\hline
\textnormal{No. of Girders} & 5 \\[6pt]
\hline
\textnormal{Girder Spacing (m)} & 2.28 m \\[6pt]
\hline
\textnormal{Deck Overhang Width (m)} & 1.14 m \\[6pt]
\hline
\textnormal{Deck Thickness (mm)} & 250.0 mm \\[6pt]
\hline
\textnormal{Footpath Width (m)} & 0.0 m (IRC 5 Cl. 104.3.6 min: 1.5 m) \\[6pt]
\hline
\textnormal{No. of Traffic Lanes} &  (per IRC 5 Cl. 104.3.1) \\[6pt]
\hline
\end{longtable}

\vspace{0.8em}

\begin{longtable}{|L{5.5cm}|p{10.0cm}|}
\caption{\textbf{Components Details}}
\hline
\textbf{parameter} & \textbf{value} \\
\hline
\textnormal{Crash Barrier Type} & IRC 5 - RCC Crash Barrier \\[6pt]
\hline
\textnormal{Crash Barrier Load (kN/m)} & 7.557 \\[6pt]
\hline

\textnormal{Railing Type} & IRC 5 - RCC Railing \\[6pt]
\hline
\textnormal{Railing Load (kN/m)} & 1.4715 \\[6pt]
\hline
\textnormal{Wearing Course Material} & Concrete \\[6pt]
\hline
\textnormal{Wearing Course Thickness (mm)} & 50.0 mm \\[6pt]
\hline
\end{longtable}


\newpage

\vspace{0.4em}
\noindent
            
\vspace{4pt}
\begin{longtable}{|L{2.2cm}|L{1.8cm}|p{3.8cm}|p{3.8cm}|p{3.8cm}|}
\caption{\textbf{Girder General Information}}
\hline
\textbf{Girder} & \textbf{Member ID} & \textbf{Design Mode} & \textbf{Girder Type} & \textbf{Girder Symmetry} \\[6pt]
\hline
G1 & G1M1 & Optimized & Welded & Girder Symmetric \\[8pt]
\hline
G2 & G2M1 & Optimized & Welded & Girder Symmetric \\[8pt]
\hline
G3 & G3M1 & Optimized & Welded & Girder Symmetric \\[8pt]
\hline
G4 & G4M1 & Optimized & Welded & Girder Symmetric \\[8pt]
\hline
G5 & G5M1 & Optimized & Welded & Girder Symmetric \\[8pt]
\hline
\end{longtable}

\vspace{0.6em}

\vspace{4pt}
\begin{longtable}{|L{1.8cm}|L{2.3cm}|L{1.8cm}|p{4.8cm}|p{4.8cm}|}
\caption{\textbf{Girder Section Dimensions}}
\hline
\textbf{Girder} & \textbf{Total Depth, D (mm)} & \textbf{Web, $t_w$ (mm)} & \textbf{Top Flange (b\textsubscript{tf}, t\textsubscript{tf}) mm} & \textbf{Bottom Flange (b\textsubscript{bf}, t\textsubscript{bf}) mm} \\[6pt]
\hline
G1 & 1670 mm & All mm & 510 mm, All mm & 510 mm, All mm \\[8pt]
\hline
G2 & 1670 mm & All mm & 510 mm, All mm & 510 mm, All mm \\[8pt]
\hline
G3 & 1670 mm & All mm & 510 mm, All mm & 510 mm, All mm \\[8pt]
\hline
G4 & 1670 mm & All mm & 510 mm, All mm & 510 mm, All mm \\[8pt]
\hline
G5 & 1670 mm & All mm & 510 mm, All mm & 510 mm, All mm \\[8pt]
\hline
\end{longtable}

\vspace{0.6em}

\vspace{4pt}
\begin{longtable}{|L{1.4cm}|p{2.2cm}|p{2.2cm}|p{3.0cm}|p{2.4cm}|p{2.2cm}|}
\caption{\textbf{Girder Restraint and Stiffener Details}}\\
\hline
\textbf{Girder} & \textbf{Torsional / Warping Restraint} & \textbf{Web Philosophy} & \textbf{Intermediate Stiffeners} & \textbf{Longitudinal Stiffeners} & \textbf{Bearing Stiffener} \\[6pt]
\hline
\endfirsthead
\endhead

G1 & Fully Restrained, Both Flanges Restrained & Thin Web with ITS & Yes; Spacing: 2440 mm; Thickness: 10 mm & No & No.: 4; Spacing: 100 mm; Thickness: 10 mm \\[8pt]
\hline
G2 & Fully Restrained, Both Flanges Restrained & Thin Web with ITS & Yes; Spacing: 2440 mm; Thickness: 10 mm & No & No.: 4; Spacing: 100 mm; Thickness: 10 mm \\[8pt]
\hline
G3 & Fully Restrained, Both Flanges Restrained & Thin Web with ITS & Yes; Spacing: 2440 mm; Thickness: 10 mm & No & No.: 4; Spacing: 100 mm; Thickness: 10 mm \\[8pt]
\hline
G4 & Fully Restrained, Both Flanges Restrained & Thin Web with ITS & Yes; Spacing: 2440 mm; Thickness: 10 mm & No & No.: 4; Spacing: 100 mm; Thickness: 10 mm \\[8pt]
\hline
G5 & Fully Restrained, Both Flanges Restrained & Thin Web with ITS & Yes; Spacing: 2440 mm; Thickness: 10 mm & No & No.: 4; Spacing: 100 mm; Thickness: 10 mm \\[8pt]
\hline
\end{longtable}



\newpage

\vspace{0.4em}
\setlength{\tabcolsep}{4pt}
\setlength\LTleft{0pt}
\setlength\LTright{\fill}

\begin{longtable}{|L{2.2cm}|L{2.2cm}|L{3.0cm}|L{2.5cm}|C{1.8cm}|C{1.8cm}|}
\caption{\textbf{Member Properties: Cross Bracing Details}}
\hline
\textbf{Location} & \textbf{Member IDs} & \textbf{Type of Bracing} & \textbf{Bracing Section} & \textbf{Spacing (m)} \\
\hline
G1 to G2 & B1M1 to B1M1 & X-Bracing & IS 100 x 100 x 10 & 15.0 m \\[6pt]
\hline
G2 to G3 & B2M1 to B2M1 & X-Bracing & IS 100 x 100 x 10 & 15.0 m \\[6pt]
\hline
G3 to G4 & B3M1 to B3M1 & X-Bracing & IS 100 x 100 x 10 & 15.0 m \\[6pt]
\hline
G4 to G5 & B4M1 to B4M1 & X-Bracing & IS 100 x 100 x 10 & 15.0 m \\[6pt]
\hline
\end{longtable}

\vspace{0.4em}
\noindent
\setlength{\tabcolsep}{4pt}
\setlength\LTleft{0pt}
\setlength\LTright{\fill}

\begin{longtable}{|L{2.2cm}|L{2.2cm}|L{3.0cm}|L{2.5cm}|C{1.8cm}|C{1.8cm}|}
\caption{\textbf{Member Properties: End Diaphragm Details}}
\hline
\textbf{Location} & \textbf{Member IDs} & \textbf{Type of Bracing} & \textbf{Bracing Section} \\
\hline
G1 to G2 & E1M1 to E1M2 & Cross Bracing & IS 100 x 100 x 10 \\[6pt]
\hline
G2 to G3 & E2M1 to E2M2 & Cross Bracing & IS 100 x 100 x 10 \\[6pt]
\hline
G3 to G4 & E3M1 to E3M2 & Cross Bracing & IS 100 x 100 x 10 \\[6pt]
\hline
G4 to G5 & E4M1 to E4M2 & Cross Bracing & IS 100 x 100 x 10 \\[6pt]
\hline
\end{longtable}



\label{subsec:shear-connectors}

\vspace{2.2em}

\vspace{0.4em}
\begin{longtable}{|L{5.5cm}|p{10.0cm}|}
\caption{\textbf{Shear Connector Details}}
\hline
\textbf{parameter} & \textbf{value} \\
\hline
\textnormal{Stud Diameter (mm)} & 20.0 mm \\[6pt]
\hline
\textnormal{Stud Height (mm)} & 100.0 mm \\[6pt]
\hline
\textnormal{Stud $f_y$ (MPa)} & 385.0 MPa \\[6pt]
\hline
\textnormal{Stud $f_u$ (MPa)} & 495.0 MPa \\[6pt]
\hline
\textnormal{No. of Studs per Section} & 2 \\[6pt]
\hline
\end{longtable}



\label{subsec:safety-factors}

\vspace{2.2em}

\vspace{0.3em}
\noindent\textit{Note: All values are per IRC 22 Table 1 unless user-modified.}

\vspace{0.4em}
\begin{longtable}{|L{5.5cm}|p{10.0cm}|}
\caption{\textbf{Partial Safety Factors}}
\hline
\textbf{parameter} & \textbf{value} \\
\hline
\textnormal{$\gamma_{M0}$ (Yielding / Buckling)} & 1.10 \\[6pt]
\hline
\textnormal{$\gamma_{M1}$ (Ultimate Stress)} & 1.25 \\[6pt]
\hline
\textnormal{$\gamma_C$ (Concrete, Basic)} & 1.50 \\[6pt]
\hline
\textnormal{$\gamma_s$ (Reinforcement)} & 1.15 \\[6pt]
\hline
\textnormal{$\gamma_v$ (Shear Connectors)} & 1.25 \\[6pt]
\hline
\textnormal{$\gamma_{fft}$ (Fatigue Load)} & 1.00 \\[6pt]
\hline
\textnormal{$\gamma_{Mft}$ (Fatigue Strength)} & 1.35 \\[6pt]
\hline
\end{longtable}


\chapter{Material Take-off \& Quantity Summary}
\label{ch:material-takeoff}

\noindent\textbf{Table 7.1  Bill of Materials (Steel, Concrete, and Reinforcement Quantities)}

\begingroup
\setlength{\tabcolsep}{3.5pt}
\begin{longtable}{|C{1.0cm}|L{3.8cm}|C{2.6cm}|C{1.8cm}|C{1.8cm}|C{1.8cm}|C{1.8cm}|}
\hline
\textbf{S.N.} & \textbf{Item Description} & \textbf{Volume} & \textbf{Quantity} & \textbf{Total Volume} & \textbf{Weight (MT)} & \textbf{Total Weight (MT)} \\
\hline
1 & Structural Steel (IS 2062) for Girders & $0.03492\text{ m}^2 \times 30.00\text{ m} = 1.04756\text{ m}^3$ & 5 & 5.24 & 8.22 & 41.12 \\
\hline
2(a) & Cross Bracing - Top Chord & $0.00024\text{ m}^2 \times 2.28\text{ m} = 0.00054\text{ m}^3$ & 36 & 0.02 & 0.0042 & 0.15 \\
\hline
2(b) & Cross Bracing - Bottom Chord & $0.00024\text{ m}^2 \times 2.28\text{ m} = 0.00054\text{ m}^3$ & 36 & 0.02 & 0.0042 & 0.15 \\
\hline
2(c) & Cross Bracing - Diagonal Chord & $0.00024\text{ m}^2 \times 2.69\text{ m} = 0.00064\text{ m}^3$ & 72 & 0.05 & 0.0050 & 0.36 \\
\hline
3 & Concrete (M40) for Deck Slab & $11.40\text{ m} \times 0.25\text{ m} \times 30.00\text{ m} = 85.50\text{ m}^3$ & 1 & 85.50 & 213.75 & 213.75 \\
\hline
4 & Reinforcement Steel (Fe 500) & $0.043567\text{ m}^2 \times 30.00\text{ m} = 1.30701\text{ m}^3$ & 1 & 1.31 & 10.26 & 10.26 \\
\hline
5 & Shear Stud Connectors & $0.000314\text{ m}^2 \times 0.100\text{ m} = 0.000031\text{ m}^3$ & 870 & 0.03 & 0.000247 & 0.215 \\
\hline
6 & Crash Barrier & $0.30228\text{ m}^2 \times 30.00\text{ m} = 9.06848\text{ m}^3$ & 2 & 18.14 & 22.67 & 45.34 \\
\hline
\end{longtable}


\chapter{Standards \& Assumptions}
\label{ch:Design Standards}

This section provides references to standards used to calibrate the OsdagBridge
design modules, and notes the limitations of the current software version.

\section{Design Standards}
\label{sec:design_standards}

The following Indian Road Congress (IRC) codes and Indian Standards (IS) 
form the basis of all design calculations in this software.

\vspace{0.5cm}

\begingroup
\setlength{\tabcolsep}{3.5pt}
\begin{table}[H]
\caption{\textbf{IRC Codes}}
\begin{tabular}{|c|c|p{13cm}|}
\hline
\textbf{Code} & \textbf{Year} & \textbf{Title / Scope} \\ 
\hline
IRC 5 & 2015 & General Features of Design - carriageway widths, kerb, footpath dimensions \\ 
\hline
IRC 6 & 2017 & Loads and Load Combinations - dead load, live load, impact, wind, temperature, etc. \\ 
\hline
IRC 22 & 2015 & Composite Construction (LS) - Composite section properties, ULS/SLS design, shear connectors \\ 
\hline
IRC 24 & 2010 & Steel Road Bridges (LS) - Stiffener design, skew angle limits, diaphragm requirements \\ 
\hline
IRC 112 & 2020 & Concrete Road Bridges - Deck slab flexure, shear, crack width, reinforcement \\ 
\hline
IRC SP 114 & 2018 & Seismic Design of Road Bridges \\ 
\hline
\end{tabular}
\end{table}

\begingroup
\setlength{\tabcolsep}{3.5pt}
\begin{table}[H]
\caption{\textbf{IS Codes}}
\begin{tabular}{|c|c|p{13cm}|}
\hline
\textbf{Code} & \textbf{Year} & \textbf{Scope} \\
\hline
IS 800 & 2007 & Steel construction - tension, compression, bending, shear, LTB, stiffeners, combined checks \\
\hline
IS 456 & 2000 & Concrete - simplified stress-block for deck moment capacity \\
\hline
IS 1786 & 2008 & Reinforcement steel properties \\
\hline
IS 1893 (Part 3) & 2014 & Earthquake resistant design \\
\hline
IS 2062 & 2011 & Structural steel - yield and ultimate strength by grade \\
\hline
\end{tabular}
\end{table}

\clearpage
\section{Analysis and Design Assumptions of This Version}
\label{sec:assumptions}

\textbf{Structural Analysis}

\begin{itemize}
    \item All girders are modelled as simply supported; continuous spans are not currently supported.
    \item A 3D grillage model (OSPGrillage) is used for load distribution. Grillage members carry composite section properties after the construction stage.
    \item The transverse member forces are computed using an approximate 2D frame analogy. For irregular or skewed geometries, a 3D FEM is recommended.
    \item Fixed bearing stiffness is modelled as $k = 1{,}000{,}000 \,\text{kN/m}$ (virtually rigid); free/expansion bearing as $k = 100 \, \text{kN/m}$.
    \item Construction stage sequence analysis is approximate. Detailed staged analysis should be performed for long-term deflection checks.
\end{itemize}

\textbf{Composite Action}

\begin{itemize}
    \item Full shear connection is assumed at ULS with headed stud connectors designed per IRC 22:2015 Cl.606.
    \item Short-term composite section properties (modular ratio $n = E_s/E_{cm}$) are used for ULS checks.
    \item Long-term composite section properties (with creep-adjusted modular ratio) are used for SLS deflection and crack-width checks.
    \item Pre-composite stage: the steel girder alone resists all construction loads prior to concrete gaining strength.
\end{itemize}


\textbf{Material Properties (IRC 22:2015 Annex III)}

\begin{itemize}
\item Steel: $E_s = 200{,}000 \, \text{MPa}, \ G_s = 80{,}000 \, \text{MPa}, \ \nu = 0.30, \ \alpha = 11.7 \times 10^{-6}/^\circ\mathrm{C}$ (grade-independent).
\item Minimum structural concrete grade: M25.
\item Default reinforcement grade: Fe500.
\end{itemize}

\textbf{Partial Safety Factors (IRC 22:2015 Cl.601.4)}

\begin{itemize}
\item $\gamma_{m0}$ (steel yield, ULS) = 1.10
\item $\gamma_{m1}$ (steel ultimate, ULS) = 1.25
\item Reinforcement (ULS) = 1.15
\item Welds -- shop: 1.25; field: 1.50
\item Fatigue ($\gamma_{mft}$) = 1.35
\end{itemize}

\textbf{Loading}

\begin{itemize}
\item Dead load densities per IRC 6:2017 Cl.203: structural steel $78.5 \ \text{kN/m}^3$, concrete $25.0 \ \text{kN/m}^3$, bituminous wearing course $24.0 \ \text{kN/m}^3$.
\item Multi-lane live load reduction factors per IRC 6:2017 Cl.204.4 Table 6A: 1st lane = 1.0, 2nd lane = 0.8, 3rd lane onwards = 0.4.
\item Impact factor (dynamic load allowance) computed from span per IRC 6:2017 Cl.208.2/208.3.
\item Wind load applied as transverse, longitudinal, and vertical components per IRC 6:2017 Cl.209.3.3--209.3.5.
\end{itemize}

\textbf{Serviceability Limits}

\begin{itemize}
\item Deflection limits (IRC 22:2015 Cl.604.3.2): live load + impact $\leq L/800$; total $\leq L/600$.
\item SLS stress limits: concrete $\sigma_c \leq 0.48 f_{ck}$; rebar $\sigma_s \leq 0.80 f_{yk}$; steel $f_e \leq 0.9 f_y$.
\item Permissible crack width: $w_k \leq 0.3 \ \text{mm}$ (bridge deck, exposure class XS2/XD2 per IRC 112:2020 Cl.12.3.2).
\end{itemize}

\textbf{Fatigue}

\begin{itemize}
\item Reference fatigue life: $N_{sc} = 2 \times 10^6$ cycles (IRC 22:2015 Cl.605).
\item Constant stress range is assumed. A thickness correction factor $\mu_r$ is applied for plate thickness $> 25 \ \text{mm}$.
\end{itemize}

\textbf{Stiffener Design}

\begin{itemize}
\item Intermediate transverse stiffener and bearing stiffener design follows IS 800:2007 Cl.8.7.2/8.7.3 and IRC 24:2010 Cl.509.7.2/509.7.3.
\end{itemize}


\section{Known Limitations of This Version}
\label{sec:limitations}

\begin{itemize}
\item Substructure (piers, pile caps, foundations) and bearing design are not included.
\item Splice connection design is not implemented.
\item Skew angle $>$ 15 degrees requires independent manual analysis (IRC 24 Cl. 504.8).
\item Construction stage sequence analysis is approximate; detailed staged analysis
  should be performed for long-term deflection checks.
\item The grillage analysis assumes simply supported boundary conditions;
  continuous spans are not currently supported.
\end{itemize}


\chapter*{References}
\addcontentsline{toc}{chapter}{References}


\begin{enumerate}

\item IRC 5 (2024) --- \textit{Standard Specifications and Code of Practice for Road Bridges, Section I: General Features of Design.}

\item IRC 6 (2017) --- \textit{Standard Specifications and Code of Practice for Road Bridges, Section II: Loads and Load Combinations.}

\item IRC 22 (2014) --- \textit{Standard Specifications and Code of Practice for Road Bridges, Section VI: Composite Construction (Limit State Design).}

\item IRC 24 (2010) --- \textit{Standard Specifications and Code of Practice for Road Bridges, Section V: Steel Road Bridges (Limit State Method).}

\item IRC 112 (2011) --- \textit{Code of Practice for Concrete Road Bridges.}

\item IRC SP 114 (2018) --- \textit{Guidelines for Seismic Design of Road Bridges.}

\item IS 800 (2007) --- \textit{Indian Standard: General Construction in Steel --- Code of Practice.}

\item IS 2062 (2011) --- \textit{Hot Rolled Medium and High Tensile Structural Steel --- Specification.}

\item Subramanian, N. (2008). \textit{Design of Steel Structures.} Oxford University Press.

\item Steel-INSDAG Teaching Resource Materials. \url{https://www.steel-insdag.org}

\end{enumerate}

\end{document}